% Derived rules of GPA_\rtop used throughout the FME derivation.
\begin{tabular}{r c c c}
  \diagrow{(a)}{
      \pic (g1) at (0.500,0.500) {antipode};
      \begin{scope}[shift={(1.250,0.000)}]
      \pic (g2) at (0.500,0.500) {ge};
      \end{scope}
      \draw (g1-out-0) to[out=0,in=180] (g2-in-0);
  }{
      \pic (g1) at (0.500,0.500) {coge};
      \begin{scope}[shift={(1.250,0.000)}]
      \pic (g2) at (0.500,0.500) {antipode};
      \end{scope}
      \draw (g1-out-0) to[out=0,in=180] (g2-in-0);
  }
  \diagrow{(b) $k>0$}{
      \pic[val=k] (g1) at (0.500,0.500) {scalar};
      \begin{scope}[shift={(1.250,0.000)}]
      \pic (g2) at (0.500,0.500) {ge};
      \end{scope}
      \draw (g1-out-0) to[out=0,in=180] (g2-in-0);
  }{
      \pic (g1) at (0.500,0.500) {ge};
      \begin{scope}[shift={(1.250,0.000)}]
      \pic[val=k] (g2) at (0.500,0.500) {scalar};
      \end{scope}
      \draw (g1-out-0) to[out=0,in=180] (g2-in-0);
  }
  \diagrow{(b) $k<0$}{
      \pic[val=k] (g1) at (0.500,0.500) {scalar};
      \begin{scope}[shift={(1.250,0.000)}]
      \pic (g2) at (0.500,0.500) {ge};
      \end{scope}
      \draw (g1-out-0) to[out=0,in=180] (g2-in-0);
  }{
      \pic (g1) at (0.500,0.500) {coge};
      \begin{scope}[shift={(1.250,0.000)}]
      \pic[val=k] (g2) at (0.500,0.500) {scalar};
      \end{scope}
      \draw (g1-out-0) to[out=0,in=180] (g2-in-0);
  }
  \diagrow{(c)}{
      \pic (g1) at (0.500,1.000) {ge};
      \pic (g2) at (0.500,0.000) {ge};
      \pic (g3) at (1.750,0.500) {multiplication};
      \draw (g1-out-0) to[out=0,in=180] (g3-in-0);
      \draw (g2-out-0) to[out=0,in=180] (g3-in-1);
  }{
      \pic (g1) at (0.500,0.500) {multiplication};
      \pic (g2) at (1.750,0.500) {ge};
      \draw (g1-out-0) to[out=0,in=180] (g2-in-0);
  }
  \diagrow{(c$^{*}$)}{
      \pic (g1) at (0.500,0.500) {ge};
      \pic (g2) at (1.750,0.500) {comultiplication};
      \draw (g1-out-0) to[out=0,in=180] (g2-in-0);
  }{
      \pic (g1) at (0.500,0.500) {comultiplication};
      \pic (g2) at (1.750,1.000) {ge};
      \pic (g3) at (1.750,0.000) {ge};
      \draw (g1-out-0) to[out=0,in=180] (g2-in-0);
      \draw (g1-out-1) to[out=0,in=180] (g3-in-0);
  }
  \diagrow{(d)}{
      \pic (g1) at (0.500,0.500) {ge};
      \pic (g2) at (1.750,0.500) {discard};
      \draw (g1-out-0) to[out=0,in=180] (g2-in-0);
  }{
      \pic (g1) at (0.500,0.500) {discard};
      \coordinate (id_in_1) at (-0.250,0.500);
      \draw (id_in_1) -- (g1-in-0);
  }
  \diagrow{(d$^{*}$)}{
      \pic (g1) at (0.500,0.500) {codiscard};
      \pic (g2) at (1.750,0.500) {ge};
      \draw (g1-out-0) to[out=0,in=180] (g2-in-0);
  }{
      \pic (g1) at (0.500,0.500) {codiscard};
      \coordinate (id_out_1) at (1.250,0.500);
      \draw (g1-out-0) -- (id_out_1);
  }
  \diagrowc{(e)}{
      \pic (g1) at (0.500,1.000) {ge};
      \pic (g2) at (0.500,0.000) {ge};
      \pic (g3) at (1.750,0.500) {cocopy};
      \draw (g1-out-0) to[out=0,in=180] (g3-in-0);
      \draw (g2-out-0) to[out=0,in=180] (g3-in-1);
  }{
      \pic (g1) at (0.500,0.500) {cocopy};
      \pic (g2) at (1.750,0.500) {ge};
      \draw (g1-out-0) to[out=0,in=180] (g2-in-0);
  }
  \diagrow{(f) transitivity}{
      \pic (g1) at (0.500,0.500) {ge};
      \pic (g2) at (1.750,0.500) {ge};
      \draw (g1-out-0) to[out=0,in=180] (g2-in-0);
  }{
      \pic (g1) at (0.500,0.500) {ge};
  }
  \diagrow{(f) direction}{
      \pic (g1) at (0.500,0.500) {coge};
      \pic (g2) at (1.750,0.500) {ge};
      \draw (g1-out-0) to[out=0,in=180] (g2-in-0);
  }{
      \pic (g1) at (0.500,0.500) {discard};
      \pic (g2) at (1.750,0.500) {codiscard};
      \coordinate (id_in_1) at (-0.250,0.500);
      \draw (id_in_1) -- (g1-in-0);
      \coordinate (id_out_1) at (2.500,0.500);
      \draw (g2-out-0) -- (id_out_1);
  }
  \diagrow{(g)}{
      \pic (g1) at (0.500,1.000) {ge};
      \draw (g1-out-0) -- (1.600,1.000) to[out=0,in=0] (1.600,0.000) -- (0.000,0.000);
  }{
      \pic (g1) at (0.500,0.000) {coge};
      \draw (g1-out-0) -- (1.600,0.000) to[out=0,in=0] (1.600,1.000) -- (0.000,1.000);
  }
  \diagrow{(g) $k\neq0$}{
      \pic[val=k] (g1) at (0.500,1.000) {scalar};
      \draw (g1-out-0) -- (1.600,1.000) to[out=0,in=0] (1.600,0.000) -- (0.000,0.000);
  }{
      \pic[val=k^{-1}] (g1) at (0.500,0.000) {scalar};
      \draw (g1-out-0) -- (1.600,0.000) to[out=0,in=0] (1.600,1.000) -- (0.000,1.000);
  }
\end{tabular}
